\documentclass[12pt]{article}
\usepackage[utf8]{inputenc}
\usepackage[T1]{fontenc}
\usepackage[margin=1in]{geometry}
\usepackage{tabularx}
\usepackage{booktabs}
\usepackage{hyperref}
\usepackage{float}
\usepackage{xcolor}
\usepackage{array}
\usepackage{parskip}
\usepackage{listings}
\usepackage{mdframed}
\usepackage{enumitem}
\usepackage{amsmath}

\definecolor{headerblue}{HTML}{1F5C99}
\definecolor{codeblue}{HTML}{E8F0FA}
\definecolor{warnorange}{HTML}{FFF3CD}
\definecolor{warnborder}{HTML}{C47A00}
\definecolor{notebg}{HTML}{EAF4EA}
\definecolor{noteborder}{HTML}{2E7D32}
\definecolor{darkgray}{HTML}{333333}

\newcommand{\fn}[1]{{\footnotesize\ttfamily #1}}
\newcommand{\sectionrule}{%
  \vspace{2pt}{\color{headerblue}\rule{\linewidth}{1.5pt}}\vspace{6pt}}

\newcolumntype{L}[1]{>{\raggedright\arraybackslash}p{#1}}
\newcolumntype{C}[1]{>{\centering\arraybackslash}p{#1}}

\lstset{
  language=C,
  basicstyle=\footnotesize\ttfamily,
  keywordstyle=\color{headerblue}\bfseries,
  commentstyle=\color{gray}\itshape,
  backgroundcolor=\color{codeblue},
  frame=single, framerule=0.5pt,
  rulecolor=\color{headerblue!40},
  numbers=left, numberstyle=\tiny\color{gray}, numbersep=8pt,
  breaklines=true, showstringspaces=false, tabsize=4, captionpos=b,
}

\newmdenv[backgroundcolor=warnorange, linecolor=warnborder, linewidth=1.2pt,
  innerleftmargin=10pt, innerrightmargin=10pt,
  innertopmargin=6pt, innerbottommargin=6pt]{warnbox}

\newmdenv[backgroundcolor=notebg, linecolor=noteborder, linewidth=1.2pt,
  innerleftmargin=10pt, innerrightmargin=10pt,
  innertopmargin=6pt, innerbottommargin=6pt]{notebox}

\hypersetup{colorlinks=true, linkcolor=headerblue, urlcolor=headerblue,
            pdftitle={MOD-02 Flame Detection SW -- Detailed Technical Reference}}

% ════════════════════════════════════════════════════════════════════════════
\begin{document}

% ── COVER PAGE ───────────────────────────────────────────────────────────────
\begin{titlepage}
  \centering
  \vspace*{0.5in}
  {\Huge\bfseries CSE 396 --- Computer Engineering Project}\\[0.4in]
  {\huge\bfseries MOD-02: Flame Detection Software}\\[0.2in]
  {\Large Detailed Technical Reference}\\[0.6in]
  {\large\bfseries Responsible Authors:}\\[0.2in]
  \begin{tabular}{ll}
    \textbf{Name}              & \textbf{Student ID} \\ \midrule
    Mehmet Alp ATAY   & 220104004020 \\
    Adil Mert ERGORUN & 220104004048 \\
  \end{tabular}\\[0.8in]
  {\large Autonomous Fire Detection and Suppression Rover\\Group 2}\\[0.3in]
  {\large Version 0.1 \quad--\quad April 2026}
\end{titlepage}

\tableofcontents
\newpage

% ════════════════════════════════════════════════════════════════════════════
\section{Module Overview}
\sectionrule

MOD-02 is the software layer that reads raw ADC values from MOD-01 (Flame Sensing HW) and produces structured, semantically meaningful flame detection results consumed by the Autonomy FSM (MOD-07). It is the only module that knows the \emph{meaning} of IR sensor readings: it applies thresholding, determines which sensor is dominant, and assigns a directional label.

\begin{table}[H]
\centering\small
\begin{tabularx}{\textwidth}{L{2.8cm} L{3cm} X}
\toprule
\textbf{Direction} & \textbf{Module} & \textbf{Interface} \\
\midrule
Upstream (input)    & MOD-01 Flame HW  & \fn{flame\_hw\_read\_raw(uint16\_t *buffer)} -- populates \fn{raw[3]} \\
Downstream (output) & MOD-07 Autonomy  & \fn{flame\_data\_t} via \fn{flame\_read()} or callback \\
\bottomrule
\end{tabularx}
\caption{MOD-02 inter-module connections}
\end{table}

\textbf{Header file:} \fn{CSE\_396\_HW2/Header\ files/MOD\ 02/flame\_detection.h}

% ════════════════════════════════════════════════════════════════════════════
\section{Constants \& Macros}
\sectionrule

\begin{table}[H]
\centering\small
\begin{tabularx}{\textwidth}{L{5cm} L{2cm} X}
\toprule
\textbf{Macro} & \textbf{Value} & \textbf{Rationale} \\
\midrule
\fn{FLAME\_SENSOR\_COUNT}        & \fn{3}    & Three KY-026 IR sensors mounted left, center, right. Mirrors \fn{FLAME\_HW\_SENSOR\_COUNT} from MOD-01. \\
\fn{FLAME\_THRESHOLD}            & \fn{500}  & ADC reading (0--1023) above which a sensor is considered to be detecting a flame. Must be validated experimentally in the target environment. \\
\fn{FLAME\_SAMPLE\_INTERVAL\_MS} & \fn{50}   & Sampling period in milliseconds (20~Hz). Chosen as a balance between responsiveness and CPU load on the ESP32. \\
\fn{FLAME\_ADC\_MAX}             & \fn{1023} & Full-scale value of the 10-bit ADC on the ESP32. Used for normalisation if needed. \\
\bottomrule
\end{tabularx}
\caption{MOD-02 constants and their rationale}
\end{table}

% ════════════════════════════════════════════════════════════════════════════
\section{Data Types}
\sectionrule

\subsection{\fn{flame\_status\_t} --- Return / Error Codes}

\begin{table}[H]
\centering\small
\begin{tabularx}{\textwidth}{L{4.5cm} C{1cm} X}
\toprule
\textbf{Enumerator} & \textbf{Value} & \textbf{Meaning \& when returned} \\
\midrule
\fn{FLAME\_OK}              & \fn{0}  & Operation completed successfully. Returned by all functions on the happy path. \\
\fn{FLAME\_ERR\_INIT}       & \fn{-1} & Hardware initialisation failed. Returned by \fn{flame\_init()} when MOD-01's \fn{flame\_hw\_init()} reports an ADC configuration error. The module must not be used until a successful re-initialisation. \\
\fn{FLAME\_ERR\_NO\_SENSOR} & \fn{-2} & One or more sensors not responding. Indicates that at least one ADC channel returns 0 or 1023 persistently, suggesting a wiring or hardware fault. \\
\fn{FLAME\_ERR\_TIMEOUT}    & \fn{-3} & Returned by \fn{flame\_read()} when no new sample has been produced yet (e.g.\ first call before the first sampling cycle completes). \\
\bottomrule
\end{tabularx}
\caption{\fn{flame\_status\_t} enumerator definitions}
\end{table}

\subsection{\fn{flame\_direction\_t} --- Detected Flame Direction}

\begin{table}[H]
\centering\small
\begin{tabularx}{\textwidth}{L{4cm} C{1cm} X}
\toprule
\textbf{Enumerator} & \textbf{Value} & \textbf{Meaning} \\
\midrule
\fn{FLAME\_DIR\_NONE}   & \fn{0} & No sensor exceeded \fn{FLAME\_THRESHOLD}; no flame detected. Rover should remain in \fn{STATE\_SEARCH}. \\
\fn{FLAME\_DIR\_LEFT}   & \fn{1} & Left sensor (\fn{raw[0]}) has the highest reading above threshold. Rover should turn left. \\
\fn{FLAME\_DIR\_CENTER} & \fn{2} & Center sensor (\fn{raw[1]}) has the highest reading. Rover is aligned; advance forward. \\
\fn{FLAME\_DIR\_RIGHT}  & \fn{3} & Right sensor (\fn{raw[2]}) has the highest reading. Rover should turn right. \\
\bottomrule
\end{tabularx}
\caption{\fn{flame\_direction\_t} enumerator definitions}
\end{table}

\begin{notebox}
\textbf{Direction logic:} The direction is the sensor with the \emph{highest} ADC reading, provided that reading exceeds \fn{FLAME\_THRESHOLD}. Tie-breaking priority: center $>$ left $>$ right. If no sensor exceeds the threshold, direction is \fn{FLAME\_DIR\_NONE} regardless of relative values.
\end{notebox}

\subsection{\fn{flame\_data\_t} --- Aggregated Sensor Reading}

This is the primary struct shared with MOD-07. It is populated on every sampling cycle.

\begin{table}[H]
\centering\small
\begin{tabularx}{\textwidth}{L{3.8cm} L{2.5cm} L{2cm} X}
\toprule
\textbf{Field} & \textbf{Type} & \textbf{Range} & \textbf{Description} \\
\midrule
\fn{raw[FLAME\_SENSOR\_COUNT]} & \fn{uint16\_t[3]} & 0--1023 & Raw 10-bit ADC readings. Index 0 = left, 1 = center, 2 = right. Updated every cycle. \\
\fn{max\_intensity}  & \fn{uint16\_t}       & 0--1023       & Highest value among \fn{raw[0..2]}. Proxy for flame proximity: higher = closer or brighter. \\
\fn{detected}        & \fn{uint8\_t}        & 0 or 1        & Non-zero if \fn{max\_intensity > FLAME\_THRESHOLD}. Primary flame-presence flag for MOD-07. \\
\fn{direction}       & \fn{flame\_direction\_t} & enum      & Dominant direction derived from \fn{raw}. Only meaningful when \fn{detected == 1}. \\
\fn{timestamp\_ms}   & \fn{uint32\_t}       & ms since boot & Time of reading via \fn{millis()}. MOD-07 should reject data older than 100~ms. \\
\bottomrule
\end{tabularx}
\caption{\fn{flame\_data\_t} field descriptions}
\end{table}

\subsection{\fn{flame\_cb\_t} --- Sampling Callback}

\begin{lstlisting}[caption={Callback typedef}]
typedef void (*flame_cb_t)(const flame_data_t *data);
\end{lstlisting}

Invoked \emph{synchronously} at the end of each 50~ms sampling cycle with a read-only pointer to the freshly populated \fn{flame\_data\_t}.

\begin{warnbox}
\textbf{Thread-safety warning:} The callback runs inside the sampling loop. It must return quickly and must \textbf{not} call \fn{flame\_read()}, \fn{flame\_stop()}, or any blocking function. Blocking will stall the loop and cause missed readings.
\end{warnbox}

% ════════════════════════════════════════════════════════════════════════════
\section{Public API}
\sectionrule

\subsection{\fn{flame\_init()}}

\begin{lstlisting}[caption={Signature}]
flame_status_t flame_init(void);
\end{lstlisting}

\begin{table}[H]\centering\small
\begin{tabularx}{\textwidth}{L{3cm} X}
\toprule
Pre-condition  & None. Safe to call at any point, including after a failed init. \\
Post-condition & On \fn{FLAME\_OK}: ADC pins configured, module ready. On error: state undefined; retry. \\
Calls          & \fn{flame\_hw\_init()} from MOD-01. \\
Returns        & \fn{FLAME\_OK} or \fn{FLAME\_ERR\_INIT}. \\
Side effects   & Configures ESP32 ADC channels. Does \textbf{not} start sampling. \\
\bottomrule
\end{tabularx}
\end{table}

\subsection{\fn{flame\_start()}}

\begin{lstlisting}[caption={Signature}]
flame_status_t flame_start(flame_cb_t cb);
\end{lstlisting}

\begin{table}[H]\centering\small
\begin{tabularx}{\textwidth}{L{3cm} X}
\toprule
Pre-condition  & \fn{flame\_init()} returned \fn{FLAME\_OK}. \\
Parameters     & \fn{cb} -- callback per sample. Pass \fn{NULL} to use pull model only. \\
Post-condition & 20~Hz sampling begins. Each cycle: reads ADC, computes direction, updates buffer, invokes \fn{cb} if set. \\
Returns        & \fn{FLAME\_OK}. \\
Side effects   & Starts repeating timer; consumes CPU proportional to sampling rate. \\
\bottomrule
\end{tabularx}
\end{table}

\subsection{\fn{flame\_read()}}

\begin{lstlisting}[caption={Signature}]
flame_status_t flame_read(flame_data_t *out);
\end{lstlisting}

\begin{table}[H]\centering\small
\begin{tabularx}{\textwidth}{L{3cm} X}
\toprule
Pre-condition  & \fn{flame\_start()} called. \fn{out} points to valid \fn{flame\_data\_t} memory. \\
Post-condition & \fn{*out} populated with the latest sample. \\
Returns        & \fn{FLAME\_OK} if data available; \fn{FLAME\_ERR\_TIMEOUT} if no sample yet. \\
Thread safety  & Safe to call from main loop while sampling runs. Single-core copy, no locking needed. \\
\bottomrule
\end{tabularx}
\end{table}

\subsection{\fn{flame\_get\_direction()}}

\begin{lstlisting}[caption={Signature}]
flame_direction_t flame_get_direction(void);
\end{lstlisting}

\begin{table}[H]\centering\small
\begin{tabularx}{\textwidth}{L{3cm} X}
\toprule
Pre-condition  & At least one sampling cycle completed. \\
Returns        & Current dominant direction. Returns \fn{FLAME\_DIR\_NONE} if no sample yet. \\
Use case       & Lightweight FSM polling when only direction is needed. \\
\bottomrule
\end{tabularx}
\end{table}

\subsection{\fn{flame\_stop()}}

\begin{lstlisting}[caption={Signature}]
void flame_stop(void);
\end{lstlisting}

\begin{table}[H]\centering\small
\begin{tabularx}{\textwidth}{L{3cm} X}
\toprule
Pre-condition  & None. Safe to call even if not started. \\
Post-condition & Sampling halted. Sensors powered down. Last sample retained in buffer. \\
Side effects   & Pending callback cancelled. Must call \fn{flame\_init()} + \fn{flame\_start()} to resume. \\
\bottomrule
\end{tabularx}
\end{table}

% ════════════════════════════════════════════════════════════════════════════
\section{Timing \& Sampling Behaviour}
\sectionrule

\begin{table}[H]\centering\small
\begin{tabularx}{\textwidth}{L{5cm} X}
\toprule
\textbf{Parameter} & \textbf{Value / Notes} \\
\midrule
Sampling interval    & 50~ms (20~Hz) -- \fn{FLAME\_SAMPLE\_INTERVAL\_MS} \\
ADC conversion time  & $\approx$~5~\textmu s per channel on ESP32; three channels $\approx$~15~\textmu s total, negligible vs.\ interval. \\
Callback deadline    & Callback must return within $\ll$~50~ms to avoid missing the next cycle. \\
Data staleness limit & MOD-07 should reject \fn{flame\_data\_t} with \fn{timestamp\_ms} older than 100~ms (2 missed cycles). \\
\bottomrule
\end{tabularx}
\caption{MOD-02 timing parameters}
\end{table}

% ════════════════════════════════════════════════════════════════════════════
\section{Inter-Module Protocol: MOD-01 $\rightarrow$ MOD-02 $\rightarrow$ MOD-07}
\sectionrule

\begin{enumerate}[leftmargin=*, label=\textbf{Step \arabic*.}]
  \item \textbf{Hardware read.} MOD-02 calls \fn{flame\_hw\_read\_raw(raw)} (MOD-01). Three 10-bit ADC values written into \fn{raw[0..2]}: left, center, right.
  \item \textbf{Threshold check.} For each \fn{raw[i]}, compare to \fn{FLAME\_THRESHOLD = 500}.
  \item \textbf{Max intensity.} \fn{max\_intensity = max(raw[0], raw[1], raw[2])}.
  \item \textbf{Detection flag.} \fn{detected = (max\_intensity > FLAME\_THRESHOLD) ? 1 : 0}.
  \item \textbf{Direction.} Index with highest \fn{raw[i]} exceeding threshold $\rightarrow$ \fn{0=LEFT, 1=CENTER, 2=RIGHT}. If none exceed threshold $\rightarrow$ \fn{FLAME\_DIR\_NONE}.
  \item \textbf{Timestamp.} \fn{timestamp\_ms = millis()}.
  \item \textbf{Publish.} Buffer updated. Callback invoked (if registered). MOD-07 calls \fn{flame\_read()} or reads callback arg.
\end{enumerate}

% ════════════════════════════════════════════════════════════════════════════
\section{Error Handling Guide}
\sectionrule

\begin{table}[H]\centering\small
\begin{tabularx}{\textwidth}{L{4.5cm} X}
\toprule
\textbf{Error code} & \textbf{Recommended action} \\
\midrule
\fn{FLAME\_ERR\_INIT}       & Wait 100~ms, retry \fn{flame\_init()}. After 3 consecutive failures, halt and signal MOD-08 system fault. \\
\fn{FLAME\_ERR\_NO\_SENSOR} & Check ADC wiring. Treat as hardware fault; do not continue operating with faulty sensor. \\
\fn{FLAME\_ERR\_TIMEOUT}    & Retry after \fn{FLAME\_SAMPLE\_INTERVAL\_MS + 10}~ms. If persistent, call \fn{flame\_stop()} then \fn{flame\_start()} to restart the loop. \\
\bottomrule
\end{tabularx}
\caption{MOD-02 error handling recommendations}
\end{table}

% ════════════════════════════════════════════════════════════════════════════
\section{Integration Example}
\sectionrule

\begin{lstlisting}[caption={Full MOD-02 integration with callback and error handling}]
#include "flame_detection.h"
#include <Arduino.h>

/* Callback: invoked by MOD-02 every 50 ms */
void on_flame_sample(const flame_data_t *data) {
    if (!data->detected) return;

    switch (data->direction) {
        case FLAME_DIR_LEFT:
            Serial.printf("[%lu ms] Flame LEFT  | intensity=%u\n",
                          data->timestamp_ms, data->max_intensity);
            break;
        case FLAME_DIR_CENTER:
            Serial.printf("[%lu ms] Flame AHEAD | intensity=%u\n",
                          data->timestamp_ms, data->max_intensity);
            break;
        case FLAME_DIR_RIGHT:
            Serial.printf("[%lu ms] Flame RIGHT | intensity=%u\n",
                          data->timestamp_ms, data->max_intensity);
            break;
        default: break;
    }
}

void setup() {
    Serial.begin(115200);

    /* Step 1: initialise hardware (MOD-01 ADC pins) */
    if (flame_init() != FLAME_OK) {
        Serial.println("ERROR: flame_init failed. Halting.");
        while (1);
    }

    /* Step 2: start 20 Hz sampling with callback */
    flame_start(on_flame_sample);
}

void loop() {
    /* Pull model: read latest data every 100 ms */
    flame_data_t latest;
    if (flame_read(&latest) == FLAME_OK) {
        /* Forward to MOD-07 autonomy layer:                        */
        /* autonomy_input_t in = { .flame_data = latest, ... };    */
        /* autonomy_update(&in);                                    */
    }
    delay(100);
}
\end{lstlisting}

% ════════════════════════════════════════════════════════════════════════════
\section{Known Limitations \& TODOs}
\sectionrule

\begin{itemize}
  \item \textbf{Ambient light interference.} Strong sunlight or incandescent lamps generate IR radiation indistinguishable from flame by the KY-026. Shielded housings and directional tubes are recommended. \fn{FLAME\_THRESHOLD} must be tuned in the deployment environment.
  \item \textbf{Single-axis detection only.} Left/center/right; no elevation angle. A flame significantly above or below the sensor plane may not be detected.
  \item \textbf{Winner-takes-all direction.} A weighted angular average of all \fn{raw[i]} values would provide finer directional precision than the current dominant-sensor approach.
  \item \textbf{Callback thread safety.} If FreeRTOS tasks are used, a mutex must protect the shared \fn{flame\_data\_t} buffer.
  \item \textbf{TODO:} Moving-average filter over last $N$ samples before threshold comparison to reduce noise.
  \item \textbf{TODO:} Add \fn{FLAME\_ERR\_STALE} when \fn{timestamp\_ms} has not updated within $2 \times$ \fn{FLAME\_SAMPLE\_INTERVAL\_MS}.
  \item \textbf{TODO:} Validate \fn{FLAME\_THRESHOLD = 500} experimentally with the KY-026 at typical operating distances and lighting conditions.
\end{itemize}

\vfill
{\color{headerblue}\rule{\linewidth}{1.5pt}}
\begin{center}
  \small
  Mehmet Alp ATAY (220104004020) \quad--\quad Adil Mert ERGORUN (220104004048)\\
  CSE 396 Group 2 \quad--\quad Version 0.1 \quad--\quad April 2026
\end{center}

\end{document}
